% GENERATED FILE. Edit canonical lessons, then run npm run generate:books.
% canonical-source-hash: fnv1a64:460a3842441fc370
% canonical-lessons: SA-C459-avajna, SA-C459-avamanyate, SA-C459-tiraskaroti, SA-C459-adarah, SA-C459-manah

\chapter{Respect and Contempt}
\label{ch:sa-respect-and-contempt}
\hlchaptermodality{459}

\begin{chapteropening}
\textbf{By the end of this chapter:} \emph{I can say whether someone was respected or slighted.}

Complete the last lesson of chapter 459: I can say whether someone was respected or slighted.
\end{chapteropening}

\section[\texorpdfstring{\sk{अवज्ञा} (avajñā)}{अवज्ञा (avajñā)}]{\sk{अवज्ञा} (avajñā) — contempt}

\label{lesson:SA-C459-avajna}



\begin{quote}
Before the new one: say the Sanskrit for arrogant, then the Sanskrit for modest.
\end{quote}

\subsection*{You'll want to know: \sk{अवज्ञा}}
\textbf{\sk{अवज्ञा}} — \emph{avajñā} — \textquotedblleft{}contempt\textquotedblright{}.

\begin{quote}
\sk{गुरोः} \sk{अवज्ञा} \sk{न} \sk{उचिता।}
\end{quote}

\begin{quote}
\emph{guroḥ avajñā na ucitā.}
\end{quote}

\begin{quote}
Contempt for a teacher is not proper.
\end{quote}

A feminine noun.

\begin{cousinweb}
\textbf{\sk{अवज्ञा}} (\emph{avajñā}) is built from \emph{ava}, \textquotedblleft{}down\textquotedblright{}, and the root √jñā, \textquotedblleft{}to know\textquotedblright{}, the root of \textbf{\sk{जानाति}}: to think little of; MW gives \textquotedblleft{}contempt, disesteem, disrespect\textquotedblright{} (MW).
\end{cousinweb}

\begin{grammarlens}[title={What you've built}]
One more word for describing what you lived and felt.
\end{grammarlens}

\subsection*{Your turn}
\begin{itemize}
\raggedright
  \item \emph{Say it:} \emph{avajñā}
  \item \emph{Say it:} \emph{avajñā}, once more
  \item \emph{Say it:} \emph{guroḥ avajñā na ucitā.} — Contempt for a teacher is not proper.
  \item \emph{Recall:} say the Sanskrit for arrogant, then the Sanskrit for modest, then say \emph{avajñā} again
\end{itemize}

\begin{tcolorbox}[breakable,colback=teal!4,colframe=teal!35!black,title={Before you move on}]
What does \sk{अवज्ञा} mean? (\textquotedblleft{}Contempt\textquotedblright{}.) Say it once more.
\end{tcolorbox}
\section[\texorpdfstring{\sk{अवमन्यते} (avamanyate)}{अवमन्यते (avamanyate)}]{\sk{अवमन्यते} (avamanyate) — despises}

\label{lesson:SA-C459-avamanyate}



\begin{quote}
Before the new one: say the Sanskrit for modest, then the Sanskrit for contempt.
\end{quote}

\subsection*{You'll want to know: \sk{अवमन्यते}}
\textbf{\sk{अवमन्यते}} — \emph{avamanyate} — \textquotedblleft{}despises\textquotedblright{}.

\begin{quote}
\sk{धनिकः} \sk{दरिद्रम्} \sk{अवमन्यते}, \sk{तत्} \sk{अनुचितम्।}
\end{quote}

\begin{quote}
\emph{dhanikaḥ daridram avamanyate, tat anucitam.}
\end{quote}

\begin{quote}
The rich man despises the poor man; that is wrong.
\end{quote}

Past (he …): \sk{सः} \sk{दरिद्रम्} \sk{अवमतवान्}, \emph{saḥ daridram avamatavān}.

\begin{cousinweb}
\textbf{\sk{अवमन्यते}} (\emph{avamanyate}) is built from \emph{ava}, \textquotedblleft{}down\textquotedblright{}, and the root √man, \textquotedblleft{}to think\textquotedblright{}: to think down on; MW gives \textquotedblleft{}to despise, treat contemptuously\textquotedblright{} (MW).
\end{cousinweb}

\begin{grammarlens}[title={What you've built}]
One more word for describing what you lived and felt.
\end{grammarlens}

\subsection*{Your turn}
\begin{itemize}
\raggedright
  \item \emph{Say it:} \emph{avamanyate}
  \item \emph{Say it:} \emph{avamanyate}, once more
  \item \emph{Say it:} \emph{dhanikaḥ daridram avamanyate, tat anucitam.} — The rich man despises the poor man; that is wrong.
  \item \emph{Recall:} say the Sanskrit for modest, then the Sanskrit for contempt, then say \emph{avamanyate} again
\end{itemize}

\begin{tcolorbox}[breakable,colback=teal!4,colframe=teal!35!black,title={Before you move on}]
What does \sk{अवमन्यते} mean? (\textquotedblleft{}Despises\textquotedblright{}.) Say it once more.
\end{tcolorbox}
\section[\texorpdfstring{\sk{तिरस्करोति} (tiraskaroti)}{तिरस्करोति (tiraskaroti)}]{\sk{तिरस्करोति} (tiraskaroti) — insults}

\label{lesson:SA-C459-tiraskaroti}



\begin{quote}
Before the new one: say the Sanskrit for contempt, then the Sanskrit for despises.
\end{quote}

\subsection*{You'll want to know: \sk{तिरस्करोति}}
\textbf{\sk{तिरस्करोति}} — \emph{tiraskaroti} — \textquotedblleft{}insults\textquotedblright{}.

\begin{quote}
\sk{सः} \sk{अतिथिं} \sk{तिरस्करोति}, \sk{तत्} \sk{अनुचितम्।}
\end{quote}

\begin{quote}
\emph{saḥ atithiṃ tiraskaroti, tat anucitam.}
\end{quote}

\begin{quote}
He insults the guest; that is improper.
\end{quote}

Past (he …): \sk{सः} \sk{अतिथिं} \sk{तिरस्कृतवान्}, \emph{saḥ atithiṃ tiraskṛtavān}.

\begin{cousinweb}
\textbf{\sk{तिरस्करोति}} (\emph{tiraskaroti}) is built from \emph{tiras}, \textquotedblleft{}across, aside\textquotedblright{}, and the root √kṛ, \textquotedblleft{}to do\textquotedblright{}, the root of \textbf{\sk{करोति}}: to set aside; MW gives \textquotedblleft{}to set aside\textquotedblright{} and later senses of slighting (MW).
\end{cousinweb}

\begin{grammarlens}[title={What you've built}]
One more word for describing what you lived and felt.
\end{grammarlens}

\subsection*{Your turn}
\begin{itemize}
\raggedright
  \item \emph{Say it:} \emph{tiraskaroti}
  \item \emph{Say it:} \emph{tiraskaroti}, once more
  \item \emph{Say it:} \emph{saḥ atithiṃ tiraskaroti, tat anucitam.} — He insults the guest; that is improper.
  \item \emph{Recall:} say the Sanskrit for contempt, then the Sanskrit for despises, then say \emph{tiraskaroti} again
\end{itemize}

\begin{tcolorbox}[breakable,colback=teal!4,colframe=teal!35!black,title={Before you move on}]
What does \sk{तिरस्करोति} mean? (\textquotedblleft{}Insults\textquotedblright{}.) Say it once more.
\end{tcolorbox}
\section[\texorpdfstring{\sk{आदरः} (ādaraḥ)}{आदरः (ādaraḥ)}]{\sk{आदरः} (ādaraḥ) — regard}

\label{lesson:SA-C459-adarah}



\begin{quote}
Before the new one: say the Sanskrit for despises, then the Sanskrit for insults.
\end{quote}

\subsection*{You'll want to know: \sk{आदरः}}
\textbf{\sk{आदरः}} — \emph{ādaraḥ} — \textquotedblleft{}regard\textquotedblright{}.

\begin{quote}
\sk{शिष्यः} \sk{गुरुं} \sk{प्रति} \sk{आदरं} \sk{दर्शयति।}
\end{quote}

\begin{quote}
\emph{śiṣyaḥ guruṃ prati ādaraṃ darśayati.}
\end{quote}

\begin{quote}
The disciple shows regard for his teacher.
\end{quote}

A masculine noun.

\begin{cousinweb}
\textbf{\sk{आदरः}} (\emph{ādaraḥ}) is built from \emph{ā} and the root √dṛ, \textquotedblleft{}to heed\textquotedblright{}; MW gives \textquotedblleft{}respect, regard, notice\textquotedblright{} (MW).
\end{cousinweb}

\begin{grammarlens}[title={What you've built}]
One more word for describing what you lived and felt.
\end{grammarlens}

\subsection*{Your turn}
\begin{itemize}
\raggedright
  \item \emph{Say it:} \emph{ādaraḥ}
  \item \emph{Say it:} \emph{ādaraḥ}, once more
  \item \emph{Say it:} \emph{śiṣyaḥ guruṃ prati ādaraṃ darśayati.} — The disciple shows regard for his teacher.
  \item \emph{Recall:} say the Sanskrit for despises, then the Sanskrit for insults, then say \emph{ādaraḥ} again
\end{itemize}

\begin{tcolorbox}[breakable,colback=teal!4,colframe=teal!35!black,title={Before you move on}]
What does \sk{आदरः} mean? (\textquotedblleft{}Regard\textquotedblright{}.) Say it once more.
\end{tcolorbox}
\section[\texorpdfstring{\sk{मानः} (mānaḥ)}{मानः (mānaḥ)}]{\sk{मानः} (mānaḥ) — self-respect; honour}

\label{lesson:SA-C459-manah}



\begin{quote}
Before the new one: say the Sanskrit for insults, then the Sanskrit for regard.
\end{quote}

\subsection*{You'll want to know: \sk{मानः}}
\textbf{\sk{मानः}} — \emph{mānaḥ} — \textquotedblleft{}self-respect; honour\textquotedblright{}.

\begin{quote}
\sk{मानः} \sk{धनात्} \sk{महत्तरः} \sk{अस्ति।}
\end{quote}

\begin{quote}
\emph{mānaḥ dhanāt mahattaraḥ asti.}
\end{quote}

\begin{quote}
Honour is greater than wealth.
\end{quote}

A masculine noun.

\begin{cousinweb}
\textbf{\sk{मानः}} (\emph{mānaḥ}) comes from the root √man, \textquotedblleft{}to think\textquotedblright{}: what one is thought to be worth; MW gives \textquotedblleft{}regard, respect, honour\textquotedblright{} and also \textquotedblleft{}pride\textquotedblright{} (MW).
\end{cousinweb}

\begin{grammarlens}[title={What you've built}]
One more word for describing what you lived and felt.
\end{grammarlens}

\subsection*{Your turn}
\begin{itemize}
\raggedright
  \item \emph{Say it:} \emph{mānaḥ}
  \item \emph{Say it:} \emph{mānaḥ}, once more
  \item \emph{Say it:} \emph{mānaḥ dhanāt mahattaraḥ asti.} — Honour is greater than wealth.
  \item \emph{Recall:} say the Sanskrit for insults, then the Sanskrit for regard, then say \emph{mānaḥ} again
\end{itemize}

\begin{tcolorbox}[breakable,colback=teal!4,colframe=teal!35!black,title={Before you move on}]
What does \sk{मानः} mean? (\textquotedblleft{}Self-respect; honour\textquotedblright{}.) Say it once more.
\end{tcolorbox}
